\paragraph{Statement E.17 (a fixed-ball extension of the countable diagonal).}
Let $E$ be a countably infinite nonabelian simple group, and suppose
\[
 N(g):=\sup_{a\ne1}\rho_a(g)<\infty\qquad(g\in E),
\]
where $\rho_a$ is conjugacy length with respect to $a^{\pm1}$.
Assume there are $d\ne1$ and an integer $L\geq0$ such that, writing
$D=N(d)$, every conjugate of every $c\ne1$ with $N(c)\leq2D$
can be reached by a conjugator $v$ satisfying $N(v)\leq L$.
Then the conclusion of the killing-block lemma A.27 holds for $E$:
for any actual aspherical CW complex $B$ of dimension at most two,
$K\leq H=\pi_1B$, and epimorphism $q:K\twoheadrightarrow E$, there
is an actual aspherical extension $B\subseteq Y$ of dimension at
most two such that
\[
 H\hookrightarrow P=\pi_1Y,\qquad
 H_*(B;\mathbb Z)\cong H_*(Y;\mathbb Z),\qquad
 \langle\!\langle H\rangle\!\rangle_P=P,
\]
and no homomorphism $P\to E$ restricts to $q$ on $K$.
Only countably many cells are attached.

The hypothesis is weaker than requiring one transport bound on
all of $E$: it concerns just the ball of radius $2N(d)$.
In particular it follows from the bounded-conjugator condition
stated as Assumption~4.1 in
\cite{EGismatullinMajcherZiegler2024}, applied to the constant
family $(E,N)$. No ultraproduct theorem from that paper is used.

\emph{Proof.}
The function $N$ is integer-valued, symmetric, conjugation-invariant,
and subadditive: each of these assertions follows from the
corresponding property of every $\rho_a$ by taking the supremum.
Put $C=L+D$, $m=2C$, and
$B_C=\{g\in E:N(g)\leq C\}$. This ball contains $1$ and $d$.
For $a\in E$ define
\[
 W_a(u_1,v_1,\ldots,u_m,v_m)
       =\prod_{i=1}^{m}v_i[a,u_i]v_i^{-1}.
 \tag{E.17.1}
\]
We claim that for every $a\ne1$,
\[
             B_{2C}\subseteq W_a(B_C^{\,2m}).
 \tag{E.17.2}
\]
Choose a conjugate $u$ of $d$ with $c=[a,u]\ne1$.
Such a conjugate exists because the conjugates of $d$ normally
generate the simple group $E$; if all commuted with $a$, then
$a$ would be central. We have $N(u)=D$ and
$N(c)\leq2D$. Each $g\in B_{2C}$ is a product of at most
$\rho_c(g)\leq N(g)\leq2C=m$ conjugates of $c^{\pm1}$.
For every factor, the assumed transport bound replaces its
conjugator by one with $N$ at most $L$. Positive factors now
have the form in (E.17.1), with $N(u_i)=D$ and $N(v_i)\leq L$.
For a negative factor use
\[
                  c^{-1}=u[a,u^{-1}]u^{-1}.
\]
Its new conjugator has norm at most $L+D=C$, and its commutator
variable still has norm $D$. Set unused variables equal to $1$.
All variables lie in $B_C$, proving (E.17.2). Consequently,
for every $a\ne1$ and $b\in B_C$,
\[
                  B_C\subseteq W_a(B_C^{\,2m})b,
 \tag{E.17.3}
\]
since $N(gb^{-1})\leq2C$ whenever $g,b\in B_C$.
Notice that no bound on $N(a)$ is required.

Enumerate $E$ without repetition as $e_1,e_2,\ldots$ and
choose $h\ne1$ in $E$. Use one variable $x_s$ at each vertex
of the rooted $2m$-ary tree, with root $o$. We shall choose
coefficients $b_s\in B_C$, starting with $b_o=1$, and impose
\[
 \begin{aligned}
 x_o&=W_h(x_1,\ldots,x_{2m})b_o,\\
 x_s&=W_{x_oe_n^{-1}}(x_{s1},\ldots,x_{s,2m})b_s
            &&(|s|=n\geq1).
 \end{aligned}
 \tag{E.17.4}
\]
After coefficients at depths less than $n$ have been fixed,
substitute $x_o=e_n$ into those earlier right sides. The
nonroot coefficients $e_ne_j^{-1}$, $j<n$, are nonidentity.
The resulting map $T_n$ from depth-$n$ values to the computed
root value satisfies
\[
                  B_C\subseteq T_n(B_C^{(2m)^n}).
 \tag{E.17.5}
\]
Indeed, lift any prescribed root value in $B_C$ backwards
through the finite tree, using (E.17.3) at each vertex. The
lists of child variables are disjoint, so the choices do not
interfere and all chosen intermediate values remain in $B_C$.
As $B_C$ has at least two elements, choose a value in $B_C$
different from $e_n$ and take a corresponding tuple
$(b_s)_{|s|=n}\in B_C^{(2m)^n}$ in (E.17.5).

A solution in $E$ would have $x_o=e_n$ for some $n$.
At depth $n$ every $W$ has coefficient $1$, so (E.17.4)
forces $x_s=b_s$ there. Earlier equations then force
$x_o=T_n((b_s))\ne e_n$, a contradiction. Thus this countable
system has no solution in $E$.

Lift every chosen coefficient through $q$ into $K\leq H$
and attach the displayed variable loops and relator two-cells
to $B$. Each $W$ has exponent zero in its coefficient and
in every child variable, so the relative cellular boundary
sends each relator precisely to its own variable. It is the
identity on the countable free abelian direct sums, proving
the integral homology isomorphism.

For the subcomplex containing the cells down to depth $n$,
eliminate the nonroot generator--relator pairs in descending
order of depth. Only free frontier loops, the root loop,
and one root relator remain over $B$. That last relator has
root exponent one, since the outer word $W_h$ has exponent
zero in every child. The coefficient group is the torsion-free
free product of $H$ with the remaining free group.
Forester--Rourke's unimodular adjunction theorem, in the
form audited in E.7, gives an aspherical finite-depth
subcomplex and an injective coefficient map. This invocation
uses \cite[Main Theorem]{EForesterRourkeDiagrams}; it does not
assume that an arbitrary identity-exponent system is aspherical.
Every sphere or nullhomotopy in the increasing CW union is
contained in a finite-depth subcomplex, so $Y$ is aspherical
and $H\to\pi_1Y$ is injective.

Finally, killing $H$ makes the root coefficient and $b_o$
trivial, so $x_o=1$. All other coefficients then become
trivial, and (E.17.4) makes every variable $1$. Hence $H$
normally generates $P$. A homomorphism $P\to E$ extending
$q$ would supply the impossible solution above.
\hfill$\square$

The same enumeration and $\omega_1$-stage bookkeeping as in
A.28 can therefore exclude these additional targets while
preserving the actual two-dimensional model and its homology.
This does not settle E2. A remaining countable simple target
with pointwise finite $N$ must now fail the stated transport
bound for every $d\ne1$: for every $L$ there must be some
$c\ne1$ with $N(c)\leq2N(d)$ and a conjugate of $c$ that
has no conjugator of norm at most $L$. The elements $c$ may
vary with $L$; no fixed-conjugacy-class assertion is inferred.
